%%%% utfpr-article.tex, 2024/10/05
%%%% Copyright (C) 2016-2024 Luiz E. M. Lima (luizeduardomlima@gmail.com)
%%
%% This work may be distributed and/or modified under the conditions of the
%% LaTeX Project Public License, either version 1.3 of this license or (at your
%% option) any later version.
%% The latest version of this license is in
%%   https://www.latex-project.org/lppl.txt
%% and version 1.3 or later is part of all distributions of LaTeX version
%% 2005/12/01 or later.
%%
%% This work has the LPPL maintenance status `maintained'.
%%
%% The Current Maintainer of this work is Luiz E. M. Lima.
%%
%% This work consists of all files listed in readme.html.
%%
%% A brief description of this work is in readme.txt.

%% Detecção e aviso sobre comandos obsoletos
% \RequirePackage[l2tabu, orthodox]{nag}

%% Classe de documento
\documentclass[%% Opções (^ = padrão; > = para pacotes):
a4paper,%% Tamanho de papel: a4paper, letterpaper (^), etc.
12pt,%% Tamanho de fonte: 10pt (^), 11pt ou 12pt
fleqn,%% Alinhamento de equações à esquerda (centralizado por padrão)
%   draft,%% Aparência de documento (>): draft ou final (^)
english,%% Idioma secundário (penúltimo) (>)
brazilian,%% Idioma primário (último) (>)
]{article}

%% Pacotes utilizados
\usepackage{bookmark}
\usepackage{tikz}
\usepackage{pgfplots}
\usepackage{mathrsfs}
\usepackage{adjustbox}
\usepackage{graphicx}
\usepackage[%% Opções (^ = padrão; ¹ = Leiaute A; ² = Leiaute B):
  Layout    = B,%% Leiaute (margens): A (esquerda menor + recuo) ou B (iguais) (^)
%   Font      = Times,%% Fonte principal: Arial (²), Calibri (¹), Times ou LaTeX(*)
%   TextColor = 575757,%% Cor (Modelo HTML) texto (preto ^): cinzento
%   EmphColor = 008080,%% Cor (Modelo HTML) ênfase (preto ^): verde azulado
%   LinkColor = ,%% Cor (Modelo HTML) hyperlink (azul escuro ^): nenhum modelo (cor texto)
%   Rule      = Both,%% Linha (cabeçalho e rodapé): None (²), Header, Footer ou Both (¹)
%   AffilIcon = Off,%% Ícones de e-mail, Lattes, ORCID e Instituição: On (^) ou Off
  DocLic    = None,%% Licença do documento: CC (Creative Commons) (^) ou None
%   CCLic     = BY-NC-ND,%% Licença CC: BY (^), BY-SA, BY-ND, BY-NC, BY-NC-SA ou BY-NC-ND
%   PageNum   = Off,%% Numeração de páginas: On (^) ou Off
%   SectNum   = Off,%% Numeração de seções (\section): On (²) ou Off (¹)
%   UnnumSect = Center,%% Alinhamento de seções (\section*): Center (¹) ou Left (²)
%   Caption   = Left,%% Alinhamento de legendas: Center (^) ou Left
%   Source    = Left,%% Alinhamento de fontes (autoria): Center (^) ou Left
%   ABNTCit   = NSB,%% Citação ABNT: AAY (Nome, Ano) (^), NRB (1) ou NSB [1]
%   DOIIcon   = On,%% Ícone de DOI em referências: On ou Off (^)
%   URLIcon   = On,%% Ícone de URL em referências: On ou Off (^)
%   Index     = On,%% Índice Remissivo (makeindex): On ou Off (^)
]{utfpr_article}


%%% Título: {Português}; {English}
\Title{%
  \ifbool{@LayoutA}{}{%
  Projeto de Processamento Digital de Imagens
  }%
}{%
  \ifbool{@LayoutA}{}{%
  }%
}

%%% Autor(es) (de 1 a 8, Leiaute A; de 1 a 16, Leiaute B): {Número}; {Dados}
% \Author{1}{%% Bolsista ou voluntário(a) principal (primeiro)
%   Name        = {Fábio Bays de Araujo},%
%   Email       = {fabiobays@alunos.utfpr.edu.br},%
%   Affiliation = {\UTFPRName, Toledo, Paraná, Brasil},%
% }
\Author{1}{%% Bolsista ou voluntário(a) principal (primeiro)
  Name        = {Haynner Joner Mattos},%
  Email       = {haynnermattos@alunos.utfpr.edu.br},%
  Affiliation = {\UTFPRName, Toledo, Paraná, Brasil},%
}
% \Author{3}{%% Bolsista ou voluntário(a) principal (primeiro)
%   Name        = {Kenedy de Paula Kaefer},%
%   Email       = {kenedykaefer@alunos.utfpr.edu.br},%
%   Affiliation = {\UTFPRName, Toledo, Paraná, Brasil},%
% }

%%% Instituição (UTFPR se nenhum evento): {Dados}
\InstitutionInfo{%
  Campus     = {Toledo},%
  Department = {Engenharia de Computação},%
  DeptLogo   = {LogoCursoAlfa.png},%% Arquivo de imagem em ./Logos/
}

%% Legendas e fontes (autoria) de ilustrações e tabelas em negrito
\BoldCaption%% Texto normal por padrão

%% Processamento (makeindex) de entradas (itens) do Índice Remissivo
% \makeindex%

\newcommand{\blankline}{\vspace{1\baselineskip}}

\begin{document}

\tableofcontents

\section{Introdução e Repositórios}
O objetivo desta etapa do projeto de PDI consiste em aumentar a variabilidade e a quantidade de imagens da base original em 3 vezes, aplicando transformações de intensidade e filtragem espacial linear. O ambiente de execução e armazenamento do projeto utiliza os seguintes endereços oficiais:
\begin{itemize}
    \item \textbf{Repositório Oficial no GitHub:} \\
    \url{https://github.com/HayJM/Projeto.git}
    \item \textbf{Jupyter Notebook Executável (Google Colab):} \\
    \url{https://colab.research.google.com/drive/1P4VxhGKD7LMfsBAVPdxCGIpQoDJSre4r?usp=sharing}
\end{itemize}

\section{Formalização Matemática das Transformações}

\subsection{Transformação Logarítmica}
A transformação logarítmica é utilizada para a expansão de contraste, mapeando uma gama estreita de níveis de intensidade baixa (tons escuros) em uma gama mais ampla de níveis de saída. A expressão matemática que rege esta transformação é:
\begin{equation}
s = c \cdot \log(1 + r)
\end{equation}
Onde $r$ é a intensidade do píxel original ($r \geq 0$) e $c$ é uma constante de escala calculada dinamicamente para garantir o mapeamento do valor máximo de saída no limite de 8 bits (255):
\begin{equation}
c = \frac{255}{\log(1 + \max(f))}
\end{equation}

\subsection{Transformação Exponencial (Lei da Potência / Correção Gamma)}
A transformação exponencial permite expandir ou comprimir o contraste de forma não linear através de uma curva controlada pelo parâmetro $\gamma$ (gamma):
\begin{equation}
s = c \cdot r^{\gamma}
\end{equation}
Onde $r$ representa a intensidade do píxel original previamente normalizada para o intervalo $[0, 1]$. Para o algoritmo implementado, adotou-se o expoente de controle $\gamma = 0.5$ para expandir os tons escuros.

\subsection{Filtro da Média usando Convolução}
O filtro da média é uma operação de filtragem espacial linear de suavização (passa-baixo), que substitui o valor de cada píxel pela média aritmética das intensidades contidas em uma vizinhança retangular definida por um \textit{kernel} (máscara). A operação de convolução discreta bidimensional para uma vizinhança de tamanho $m \times n$ é dada por:
\begin{equation}
g(x, y) = \sum_{s=-a}^{a} \sum_{t=-b}^{b} w(s, t) \cdot f(x+s, y+t)
\end{equation}

Para o Filtro da Média com Kernel $5 \times 5$, todos os 25 coeficientes possuem pesos idênticos e normalizados, conforme representado matricialmente de forma completa:
\begin{equation}
w = \frac{1}{25} \begin{bmatrix}
1 & 1 & 1 & 1 & 1 \\
1 & 1 & 1 & 1 & 1 \\
1 & 1 & 1 & 1 & 1 \\
1 & 1 & 1 & 1 & 1 \\
1 & 1 & 1 & 1 & 1
\end{bmatrix}
\end{equation}

A equação simplificada aplicada a cada píxel do mapa de intensidade resulta em:
\begin{equation}
g(x, y) = \frac{1}{25} \sum_{i=-2}^{2} \sum_{j=-2}^{2} f(x+i, y+j)
\end{equation}

\section{Estrutura de Execução e Classes}
O bloco de código implementado detecta e processa dinamicamente as 10 classes da base de dados coletada:
\texttt{00\_Cherry\_Tomatoes}, \texttt{01\_Fig}, \texttt{02\_Gala\_Apple}, \texttt{03\_Pear}, \texttt{04\_Persimmon}, \texttt{05\_Pitanga}, \texttt{06\_Pomegranate}, \texttt{07\_Sicilian\_Lemon}, \texttt{08\_Strawberry} e \texttt{09\_Tangerine}.

A rotina realiza o clone automático do repositório, efetua as transformações matemáticas citadas e conserva os resultados gerando a estrutura final nomeada como \texttt{augmented\_dataset}.




\end{document}
